\chapter{Appendectomy}

\section*{What Is This Test or Procedure?}
An appendectomy is surgery to remove the appendix, usually to treat acute appendicitis or a perforated appendix.

\section*{Alternative Names}
Appendicectomy, appendix removal, laparoscopic appendectomy

\section*{Principle}
Removing the inflamed or infected appendix prevents progression to rupture, peritonitis, abscess, and sepsis.

\section*{Why This Test or Procedure Is Ordered}
It is performed when evaluation indicates appendicitis, particularly when the appendix has ruptured or the risk of progression is high. Selected uncomplicated cases may be treated initially with antibiotics under medical supervision.

\section*{How This Test or Procedure Is Performed}
Under general anesthesia, the appendix is removed through laparoscopic ports or an open incision. If rupture or an abscess is present, the surgeon may wash out the abdomen, drain infected fluid, and prescribe antibiotics.

\section*{Preparation, Risks, and Recovery}
Assessment may include examination, blood tests, and imaging. Risks include bleeding, wound infection, intra-abdominal abscess, bowel injury, adhesions, blood clots, and anesthesia complications. Uncomplicated cases often require a short hospital stay.

\section*{References}
\begin{enumerate}
  \item MedlinePlus. Appendectomy. U.S. National Library of Medicine; 2026.
  \item NHS. Appendicitis. 2024.
\end{enumerate}

\clearpage
